\section{Conclusions}
\label{sec:conclusions}

We have calculated the ground exciton and biexciton state of a
strongly oblate semi-ellipsoidal GaAs quantum dot for seven representative
geometries. A ground-orbital-matched importance transformation, a full
global-rotation-fixed biexciton estimator, and a safeguarded two-budget
pattern search were combined to limit the numerical sensitivity of the small
binding energy.

The corrected biexciton energy decreases as either semiaxis increases, but
the axial dimension controls the absolute energy much more strongly. At fixed
$a=5\rB$, increasing $c$ from $0.5\rB$ to $2\rB$ reduces $\EXX$ by a
factor of 24.4, whereas increasing $a$ from $3\rB$ to $10\rB$ at fixed
$c=\rB$ changes it by only $7.9\%$.

The central binding energy remains positive throughout the sampled range and
decreases monotonically as the dot is enlarged. It is resolved as positive by
the present rough numerical diagnostic at every geometry except $(10,1)$.
At that largest lateral size, $\Ebind=0.031(20)\,\Ry$ is only
marginally positive, and no firm sign claim is made.

Increasing either semiaxis also enhances the exciton coincidence overlap and
oscillator strength. This effect outweighs the reduction of optical
transition energy and shortens both radiative lifetimes. The calculated
exciton lifetime ranges from \SI{2.767}{ps} to \SI{1.068}{ps}; the adopted
four-channel relation gives $\tauXX=\tauX/4$.

The results therefore identify distinct geometrical controls: the short axial
dimension sets the dominant confinement-energy scale, while lateral
enlargement provides a strong route for suppressing weak biexciton binding and
enhancing radiative strength. Extensions to finite barriers, strained or
atomistic confinement, and a targeted high-resolution study of the
near-zero-binding regime are natural next steps.
